% !TITLE 侠客行
% !AUTHOR 李白
% !DYNASTY 唐
% !ORDER 10

\begin{poem}
赵客\textsuperscript{1}缦胡缨\textsuperscript{2}，吴钩\textsuperscript{3}霜雪明。
银鞍照白马，飒沓\textsuperscript{4}如流星。
十步杀一人，千里不留行\textsuperscript{5}。
事了拂衣去，深藏身与名。

闲过信陵\textsuperscript{6}饮，脱剑膝前横。
将炙啖朱亥\textsuperscript{7}，持觞劝侯嬴\textsuperscript{8}。
三杯吐然诺\textsuperscript{9}，五岳倒为轻\textsuperscript{10}。
眼花耳热后，意气素霓生\textsuperscript{11}。

救赵挥金槌\textsuperscript{12}，邯郸先震惊。
千秋二壮士，烜赫大梁城\textsuperscript{13}。
纵死侠骨香，不惭世上英。
谁能书阁下，白首太玄经\textsuperscript{14}。
\end{poem}

\begin{annotations}
\notetext{1}{赵客：赵地的壮士。古代燕赵之地（今河北一带）以出慷慨悲歌的侠客闻名。}
\notetext{2}{缦胡缨：一种粗制的胡人头巾。缦（màn），无文饰的粗帛；缨，帽带。侠客不讲究华服，穿戴朴素粗犷。}
\notetext{3}{吴钩：吴地（今江苏一带）出产的弯刀，以锋利著称。剑刃如霜雪般明亮。}
\notetext{4}{飒沓：马蹄声，形容马跑得飞快。飒沓（sà tà），疾速的样子。}
\notetext{5}{不留行：不停留、不逗留。千里奔袭，无人能阻拦。}
\notetext{6}{信陵：信陵君魏无忌，战国四公子之一，以爱养士闻名。}
\notetext{7}{朱亥：信陵君的门客，原为市井屠夫，力大无穷。信陵君窃符救赵时，朱亥用铁锤击杀魏将晋鄙。}
\notetext{8}{侯嬴：信陵君的门客，原为大梁城守门人。他为信陵君谋划了窃符救赵的计策，后因年老不能随行，自刎以谢。}
\notetext{9}{然诺：承诺。三杯酒下肚，承诺便重于一切。}
\notetext{10}{五岳倒为轻：和这份承诺相比，五岳都显得很轻。五岳，东岳泰山、西岳华山、南岳衡山、北岳恒山、中岳嵩山。}
\notetext{11}{素霓生：白虹（白色的虹霓）升起。古人认为白虹贯日是侠义之气感动天地的征兆。}
\notetext{12}{金槌：铁锤。指朱亥用铁锤击杀晋鄙、夺取兵符救赵的故事。槌（chuí）。}
\notetext{13}{烜赫大梁城：声名显赫于大梁城。烜赫（xuān hè），声威盛大；大梁，魏国都城（今开封）。}
\notetext{14}{白首太玄经：西汉学者扬雄在书阁中埋头撰写《太玄经》直到白头。李白反问：谁愿意那样皓首穷经呢？}
\end{annotations}

\begin{appreciation}
这是一首关于侠客的诗，也是李白精神世界的一幅自画像。他写侠客，其实是在写自己理想中的模样。

诗的前四句，写侠客的外形：赵地的壮士，戴着胡人的头巾，吴钩宝剑像霜雪一样明亮。银鞍白马，奔驰起来像流星一般。飒沓是马蹄声。这是一个视觉冲击力极强的形象——银、白、霜雪、流星，全是冷色调的明亮，有一种凌厉的美感。

十步杀一人，千里不留行。事了拂衣去，深藏身与名——这是侠客的精神。本领高强，却不为名利；做了大事，却悄然离去。这种功成身退的境界，是中国文化中最高贵的理想之一。李白写这四句时，大概也在想：如果我有这样的机会，我也会这样做吧。

后半部分写了信陵君窃符救赵的故事。朱亥、侯嬴都是信陵君的门客，为了报答知遇之恩，不惜以死相报。三杯吐然诺，五岳倒为轻——三杯酒下肚，承诺便重于泰山；但和这份承诺相比，五岳也显得很轻。这是何等的意气。

最后两句谁能书阁下，白首太玄经——谁能像扬雄那样，在书阁中埋头写《太玄经》，直到白头？扬雄是西汉的学者，以苦学著称。李白用反问的语气说：没人愿意那样。与其在书斋中皓首穷经，不如在天地间轰轰烈烈地活一场。纵死侠骨香，不惭世上英——即便死了，侠客的骨头也是香的，无愧于世上的英雄。

这首诗写于李白离开长安之后。他在朝廷中没有实现自己的抱负，于是转向了另一种价值的追求：侠客的自由、义气、不羁。这不是逃避，而是一种转换——当政治的路走不通时，他选择了精神的路。李白一生都在追求大鹏的高飞，但当风不再起时，他宁愿做一颗流星，在夜空中划过一道短暂而明亮的光。
\end{appreciation}
